\documentclass[tikz,border=8pt]{standalone}
\usepackage{tikz}
\usepackage{amsmath}
\usepackage{amssymb}
\usepackage{ctex}
\usetikzlibrary{calc, positioning, arrows.meta, decorations.pathreplacing, 3d}

\begin{document}

% 等距投影宏：将 3D 点 (x,y,z) 映射到 2D
% 投影公式：x2d = x - 0.45*y,  y2d = 0.30*y + z
% 方向向量:
%   a1 = (2, 0, 0)
%   a2 = (0, 2, 0)  → 投影 (-0.45*2, 0.30*2) = (-0.9, 0.6)
%   a3 = (0, 0, 2)  → 投影 (0, 2)

\newcommand{\proj}[3]{%
  % #1=x, #2=y, #3=z  → compute 2D coords
  ({(#1) - 0.45*(#2)}, {0.30*(#2) + (#3)})%
}

\begin{tikzpicture}[>=Stealth, scale=1.0]

%% ===== 以等距投影画平行六面体 =====
% 三个列向量（3D坐标）：
%   a1 = (2.5, 0, 0)
%   a2 = (0,   2.5, 0)   → 2D: (-0.45*2.5, 0.30*2.5) = (-1.125, 0.75)
%   a3 = (0,   0,   2.5) → 2D: (0, 2.5)
%
% 8 个顶点（3D → 2D 投影）：
%   O  = (0,0,0)   → (0, 0)
%   A  = O + a1    → (2.5, 0)
%   B  = O + a2    → (-1.125, 0.75)
%   C  = O + a3    → (0, 2.5)
%   AB = O+a1+a2   → (1.375, 0.75)
%   AC = O+a1+a3   → (2.5, 2.5)
%   BC = O+a2+a3   → (-1.125, 3.25)
%   ABC= O+a1+a2+a3→ (1.375, 3.25)

\coordinate (O)  at (0,    0);
\coordinate (A)  at (2.5,  0);
\coordinate (B)  at (-1.125, 0.75);
\coordinate (C)  at (0,    2.5);
\coordinate (AB) at (1.375, 0.75);
\coordinate (AC) at (2.5,  2.5);
\coordinate (BC) at (-1.125, 3.25);
\coordinate (ABC)at (1.375, 3.25);

% 后面不可见的边（虚线）
\draw[gray!50, dashed, thin] (O)  -- (B);
\draw[gray!50, dashed, thin] (O)  -- (C);
\draw[gray!50, dashed, thin] (B)  -- (AB);
\draw[gray!50, dashed, thin] (B)  -- (BC);
\draw[gray!50, dashed, thin] (C)  -- (BC);

% 填充各面（半透明）
% 正面: O-A-AC-C（z从0到a3方向，面朝前）
\fill[blue!18, opacity=0.6] (O) -- (A) -- (AC) -- (C) -- cycle;
% 顶面: C-AC-ABC-BC
\fill[blue!25, opacity=0.5] (C) -- (AC) -- (ABC) -- (BC) -- cycle;
% 右侧面: A-AB-ABC-AC
\fill[blue!12, opacity=0.5] (A) -- (AB) -- (ABC) -- (AC) -- cycle;

% 可见的边（实线）
\draw[blue!60, thin] (O)  -- (A);
\draw[blue!60, thin] (A)  -- (AC);
\draw[blue!60, thin] (AC) -- (C);
\draw[blue!60, thin] (C)  -- (O);  % 这条边被a3方向遮住，保留可见框架
\draw[blue!60, thin] (A)  -- (AB);
\draw[blue!60, thin] (AB) -- (ABC);
\draw[blue!60, thin] (ABC)-- (AC);
\draw[blue!60, thin] (ABC)-- (BC);
\draw[blue!60, thin] (BC) -- (C);

% 坐标轴（延伸至框体外）
\draw[->, thick, gray!70] (0,0)   -- (3.5, 0)
  node[right, font=\small] {$x$};
\draw[->, thick, gray!70] (0,0)   -- (-1.6, 1.1)
  node[left, font=\small] {$y$};
\draw[->, thick, gray!70] (0,0)   -- (0,   3.8)
  node[above, font=\small] {$z$};

% 原点
\fill (0,0) circle (1.8pt) node[below left, font=\scriptsize] {$O$};

% 三个列向量（粗箭头）
\draw[->, very thick, blue!80!black] (O) -- (A)
  node[below right, font=\small] {$\mathbf{a}_1\,(2,0,0)$};
\draw[->, very thick, teal!80!black] (O) -- (B)
  node[below, font=\small, xshift=-6pt] {$\mathbf{a}_2\,(0,2,0)$};
\draw[->, very thick, red!70!black] (O) -- (C)
  node[right, font=\small, xshift=2pt] {$\mathbf{a}_3\,(0,0,2)$};

% 体积标注（六面体中心）
\node[font=\normalsize\bfseries, blue!70!black] at (0.7, 1.8)
  {体积 $= |\det A|$};

%% ===== 公式框（画布外绝对坐标） =====
\node[font=\small, fill=white, draw=gray!50, thin, rounded corners=2pt,
      inner sep=7pt, align=left] at (1.5, -1.5) {
  $\det A = \det\begin{pmatrix}a_{11} & a_{12} & a_{13}\\a_{21} & a_{22} & a_{23}\\a_{31} & a_{32} & a_{33}\end{pmatrix}$
  \quad
  体积 $= |\det A|$
};

%% 整体标题
\node[font=\large\bfseries] at (1.5, 4.4)
  {行列式的几何意义：平行六面体体积};

\end{tikzpicture}
\end{document}
